\begin{table}[htbp]
\centering
\caption{Standardized Effect Sizes: Regulatory Size Thresholds and Firm Bunching}
\label{tab:sde}
\begin{threeparttable}
\begin{tabular}{lcccccc}
\toprule
Outcome & $\hat{\beta}$ & SE & SD($Y$) & SDE & SE(SDE) & Classification \\
\midrule
\multicolumn{7}{l}{\textit{Panel A: Pooled}} \\
\midrule
Excess mass, 10 emp. & -0.368 & (0.065) & 3.18 & -0.116 & (0.020) & Moderate negative \\
Density rate diff. & 0.103 & --- & 0.174 & 0.588 & --- & Large positive \\
\midrule
\multicolumn{7}{l}{\textit{Panel B: Heterogeneous (by threshold level)}} \\
\midrule
10 emp. (disclosure) & -0.368 & (0.065) & 3.18 & -0.116 & (0.020) & Moderate negative \\
50 emp. (audit+IR35) & -0.464 & (0.309) & --- & --- & --- & --- \\
\bottomrule
\end{tabular}
\begin{tablenotes}[flushleft]
\small
\item \textit{Notes:} \textbf{Country:} United Kingdom. \textbf{Research question:} Do regulatory size thresholds in the Companies Act 2006 cause private firms to distort their reported employee counts, and how large are the implied compliance costs at each threshold? \textbf{Policy mechanism:} The Companies Act 2006 defines four size categories (micro, small, medium, large) using three-dimensional thresholds on employees, turnover, and balance sheet total; crossing a boundary triggers escalating obligations including mandatory audit (\pounds5K--30K/year), IR35 off-payroll rules, Modern Slavery reporting, and gender pay gap disclosure, but only when a firm exceeds two of three criteria for two consecutive years (the ``two-of-three rule''). \textbf{Outcome definition:} Normalized excess mass ($\hat{b}$) at regulatory thresholds, measuring the fraction by which actual firm counts below the threshold exceed the polynomial-fitted counterfactual density. \textbf{Treatment:} Binary; a firm is ``treated'' by regulatory obligations when it exceeds the relevant employee threshold. \textbf{Data:} Companies House Accounts Bulk Data Product (iXBRL filings, March 2026, $N = 8{,}927$ firms across 3 filing days) and NOMIS Inter-Departmental Business Register (2010--2024, 12.4 million UK enterprises). \textbf{Method:} Polynomial bunching estimation (Chetty et al.\ 2011) with 5th-degree polynomial counterfactual and 500 bootstrap replications; aggregate density-ratio analysis comparing log-density decline rates at regulatory vs.\ non-regulatory size band boundaries. \textbf{Sample:} All UK private limited companies filing accounts with employee count data via Companies House; universe of registered enterprises from IDBR via NOMIS. SDE $= \hat{\beta} / \text{SD}(Y)$ where SD($Y$) is the pre-treatment standard deviation. Classification refers to magnitude, not statistical significance: Large ($|$SDE$| > 0.15$), Moderate ($0.05$--$0.15$), Small ($0.005$--$0.05$), Null ($< 0.005$).
\end{tablenotes}
\end{threeparttable}
\end{table}

